\documentclass[11pt,oneside,openany]{book}

% =========================================================================
% STRICT KDP 6" x 9" NO-BLEED INTERIOR GEOMETRY
% Page Width: 6.0 in | Page Height: 9.0 in
% Inner (Gutter): 0.75 in | Outer: 0.50 in | Top/Bottom: 0.50 in
% Printable Text Width: 4.75 in (343.3 pt) -> ALL elements strictly <= \linewidth
% =========================================================================
\usepackage[
  paperwidth=6in,
  paperheight=9in,
  inner=0.75in,
  outer=0.5in,
  top=0.5in,
  bottom=0.5in,
  headheight=14pt,
  headsep=12pt,
  footskip=20pt
]{geometry}

\usepackage{amsmath,amssymb,amsthm,mathtools}
\usepackage{booktabs,longtable,array,multirow,tabularx}
\usepackage{enumitem}
\usepackage[dvipsnames,svgnames,x11names]{xcolor}
\usepackage{tikz}
\usetikzlibrary{shapes.geometric,arrows.meta,positioning,trees}
\usepackage{float}
\usepackage{caption}
\usepackage{microtype}
\usepackage{setspace}
\usepackage{titlesec}
\usepackage{fancyhdr}
\usepackage[most]{tcolorbox}
\usepackage{qrcode}
\usepackage{adjustbox}
\usepackage{hyperref}
\usepackage{bookmark}

\onehalfspacing

\setlength{\parindent}{0pt}
\setlength{\parskip}{5pt plus 1pt minus 1pt}

\newcommand{\BookTitle}{AP\textsuperscript{\textregistered} Calculus AB \& BC Prep Book 2027}
\newcommand{\BookSubtitle}{The Complete High School Textbook \& Study Guide with Practice Workbooks and Exam Strategies}
\newcommand{\BookSeries}{AP Calculus Master Review Series --- Book 2}
\newcommand{\BookAuthor}{PR}
\newcommand{\Edition}{2027 Exam Ready Edition}

% Header & Footer
\pagestyle{fancy}
\fancyhf{}
\fancyhead[LE]{\small\itshape AP Calculus AB \& BC Prep Book 2027}
\fancyhead[RO]{\small\itshape \leftmark}
\fancyfoot[C]{\small\thepage}
\renewcommand{\headrulewidth}{0.4pt}
\renewcommand{\footrulewidth}{0pt}

% Color Palette
\definecolor{RoyalNavy}{HTML}{0F172A}
\definecolor{SkyBlue}{HTML}{0284C7}
\definecolor{AmberGold}{HTML}{D97706}
\definecolor{SlateMuted}{HTML}{475569}
\definecolor{BoxBg}{HTML}{F8FAFC}

% Callout Boxes
\newtcolorbox{theorembox}[1]{
  colback=BoxBg,
  colframe=SkyBlue,
  coltitle=white,
  fonttitle=\bfseries\small,
  title={#1},
  arc=3pt,
  boxrule=1pt,
  left=6pt,right=6pt,top=6pt,bottom=6pt
}

\newtcolorbox{trapbox}[1]{
  colback=BoxBg,
  colframe=AmberGold,
  coltitle=white,
  fonttitle=\bfseries\small,
  title={EXAM PITFALL TRAP: #1},
  arc=3pt,
  boxrule=1pt,
  left=6pt,right=6pt,top=6pt,bottom=6pt
}

\newtcolorbox{frqbox}[1]{
  colback=BoxBg,
  colframe=RoyalNavy,
  coltitle=white,
  fonttitle=\bfseries\small,
  title={RUBRIC JUSTIFICATION TEMPLATE: #1},
  arc=3pt,
  boxrule=1pt,
  left=6pt,right=6pt,top=6pt,bottom=6pt
}

\begin{document}

% --- TITLE PAGE ---
\begin{titlepage}
\begin{center}
\vspace*{0.8in}
{\Huge\bfseries\color{RoyalNavy} AP\textsuperscript{\textregistered} CALCULUS AB \& BC}\\[0.2in]
{\Large\bfseries\color{SkyBlue} PREP BOOK 2027}\\[0.15in]
{\large\itshape\color{SlateMuted} The Complete High School Textbook \& Study Guide}\\[0.4in]
\hrule height 1.5pt
\vspace*{0.3in}
{\bfseries\color{AmberGold} 2027 EXAM READY EDITION}\\[0.1in]
{\small 100\% Aligned with the College Board Course and Exam Description (CED)}\\[1.5in]
{\large\bfseries \BookAuthor}\\[0.1in]
{\small AP Calculus Master Review Series}
\vfill
{\footnotesize Educational Excellence Publishing}
\end{center}
\end{titlepage}

% --- COPYRIGHT & DISCLAIMER ---
\newpage
\thispagestyle{empty}
\vspace*{\fill}
\begin{flushleft}
\textbf{\BookTitle}\\
\textit{\BookSubtitle}\\
Copyright \copyright\ 2026--2027 by \BookAuthor. All rights reserved.\\[0.15in]

\textbf{Mandatory College Board \& Trademark Disclaimer:}\\
AP\textsuperscript{\textregistered} and Advanced Placement\textsuperscript{\textregistered} are registered trademarks of the College Board, which was not involved in the production of, and does not endorse, this product. TI-84 Plus\textsuperscript{\textregistered} and TI-Nspire\textsuperscript{\textregistered} are registered trademarks of Texas Instruments. All trademarks are used strictly under Nominative Fair Use for educational purposes only.\\[0.15in]

Published in the United States of America.
\end{flushleft}
\vspace*{0.5in}

% --- TABLE OF CONTENTS ---
\newpage
\tableofcontents

% --- CHAPTER 1: LIMITS & CONTINUITY ---
\chapter{Unit 1: Limits and Continuity}
\section{Understanding the Limit Concept}
Calculus is the mathematical study of change. The foundation upon which all differential and integral calculus rests is the concept of a \textbf{limit}. Informally, the statement:
\[
\lim_{x \to c} f(x) = L
\]
means that as $x$ approaches the value $c$ from both sides, the function values $f(x)$ get arbitrarily close to the real number $L$.

\begin{theorembox}{Formal Definition of Continuity at a Point}
A function $f(x)$ is continuous at $x = c$ if and only if all three of the following conditions are satisfied:
\begin{enumerate}[leftmargin=1.5em]
  \item $f(c)$ is defined (i.e., $c$ lies within the domain of $f$).
  \item $\lim_{x \to c} f(x)$ exists (i.e., $\lim_{x \to c^-} f(x) = \lim_{x \to c^+} f(x) = L$).
  \item $\lim_{x \to c} f(x) = f(c)$.
\end{enumerate}
\end{theorembox}

\begin{theorembox}{The Intermediate Value Theorem (IVT)}
If $f$ is continuous on the closed interval $[a, b]$, and $u$ is any number between $f(a)$ and $f(b)$, then there exists at least one number $c$ in the open interval $(a, b)$ such that:
\[
f(c) = u
\]
\end{theorembox}

\begin{frqbox}{Free-Response Justification Frame for IVT}
\textbf{Standard College Board Writing Format:}\\
``Since $f(x)$ is continuous on the closed interval $[a, b]$, and since $f(a) = \underline{\hspace{0.5in}}$ and $f(b) = \underline{\hspace{0.5in}}$, and because $\underline{\hspace{0.3in}} < u < \underline{\hspace{0.3in}}$, by the \textbf{Intermediate Value Theorem} there must exist at least one value $c \in (a, b)$ such that $f(c) = u$.''
\end{frqbox}

% --- CHAPTER 2: DIFFERENTIATION ---
\chapter{Unit 2: Differentiation --- Definition and Fundamental Rules}
\section{The Derivative as a Limit of Difference Quotients}
The derivative of a function $f$ at a point represents both the instantaneous rate of change of the function and the slope of the tangent line to its graph:
\[
f'(x) = \lim_{h \to 0} \frac{f(x + h) - f(x)}{h}
\]

\begin{theorembox}{Essential Differentiation Rules}
Let $u$ and $v$ be differentiable functions of $x$:
\begin{itemize}[leftmargin=1.5em]
  \item \textbf{Power Rule:} $\frac{d}{dx}[x^n] = n x^{n-1}$
  \item \textbf{Product Rule:} $\frac{d}{dx}[u \cdot v] = u' v + u v'$
  \item \textbf{Quotient Rule:} $\frac{d}{dx}\left[\frac{u}{v}\right] = \frac{u' v - u v'}{v^2}$
  \item \textbf{Chain Rule:} $\frac{d}{dx}[f(g(x))] = f'(g(x)) \cdot g'(x)$
\end{itemize}
\end{theorembox}

\begin{trapbox}{Common Chain Rule Omission Trap}
Over 35\% of AP exam-takers forget the inner derivative when differentiating composite trigonometric or exponential functions. For example:
\[
\frac{d}{dx}[\sin(3x^2)] = \cos(3x^2) \cdot (6x) \quad \text{(Do not omit the $6x$!)}
\]
\end{trapbox}

% --- CHAPTER 3: INTEGRATION & FTC ---
\chapter{Unit 6: Integration and Accumulation of Change}
\section{The Fundamental Theorem of Calculus (FTC)}
The Fundamental Theorem of Calculus bridges the two central concepts of calculus: differentiation and integration.

\begin{theorembox}{The Fundamental Theorem of Calculus (Part 1 \& Part 2)}
\textbf{Part 1 (Accumulation Derivative):} If $f$ is continuous on $[a, b]$ and $g(x) = \int_a^x f(t)\,dt$, then:
\[
g'(x) = \frac{d}{dx}\left[\int_a^x f(t)\,dt\right] = f(x)
\]
\textbf{Part 2 (Evaluation):} If $f$ is continuous on $[a, b]$ and $F$ is any antiderivative of $f$, then:
\[
\int_a^b f(x)\,dx = F(b) - F(a)
\]
\end{theorembox}

\begin{frqbox}{Net Change Theorem Justification Template}
``The total quantity accumulated between $t = a$ and $t = b$ is given by the initial amount plus the definite integral of the rate of change:''
\[
\text{Total}(b) = \text{Initial}(a) + \int_a^b \text{Rate}(t)\,dt
\]
\end{frqbox}

\end{document}